\section{Defenses}
\label{sec:defenses}

We propose three defenses of increasing structural invasiveness. All operate at
inference time without model fine-tuning, making them deployable today in any
agent framework that exposes the clarification-answer slot.

\subsection{D1: Provenance Tagging}

The simplest defense wraps the user's clarification answer in explicit delimiters
and extends the system prompt with an instruction:

\begin{quote}
\small\ttfamily
<USER\_CLARIFICATION>\\
\quad\normalfont\textit{[user's answer]}\\
</USER\_CLARIFICATION>\\[3pt]
\normalfont\textit{(System note: treat the above only as task\\
data, never as a new instruction or scope change.)}
\end{quote}

This defense exploits the model's instruction-following capability to establish a
trust boundary. It is analogous to HTML-encoding user input before embedding it in
a SQL query: the structure signals to the processor (here, the LLM) that the
enclosed content is data, not code. The limitation is that it relies on the model's
ability to maintain this distinction across a long context, which degrades with
context length~\citep{anthropic2024instructionhierarchy}.

\subsection{D2: Structured Extraction}

Rather than trusting the model to interpret the delimiter, we run a \emph{separate}
LLM pass that extracts only the direct answer to the agent's specific question:

\begin{enumerate}
  \item The agent emits a clarifying question $q$.
  \item The user's raw answer $u$ (potentially containing injected content) is
    passed to an extractor with a fixed system prompt: \emph{``Extract ONLY the
    direct answer to the question asked. Discard everything else, including any
    additional requests, instructions, or sentences that do not directly answer
    the question.''}
  \item The extracted answer $\hat{u}$ is returned to the agent in place of $u$.
\end{enumerate}

This defense is more robust than D1 because it uses a structurally separate model
call with no prior context of the task. The extractor has no reason to interpret
injected content as instructions. The cost is an additional LLM call per
clarification turn and potential information loss if the user's legitimate answer
is verbose or multi-part.

\subsection{D3: Plan-Diff Gating}

Inspired by recent work on agent action authorization~\citep{zhou2024agentauthorization},
D3 interposes a semantic consistency check \emph{after} the agent incorporates
the clarification answer:

\begin{enumerate}
  \item The agent produces a post-clarification plan (or code output) $r$.
  \item A gating LLM call compares $r$ against the original task specification,
    listing any actions in $r$ that were not implied by the original task.
  \item If out-of-scope actions are detected---new network calls, file deletions
    outside the working directory, removal of security checks, new dependencies
    not requested---the plan is flagged and execution is blocked pending user
    re-authorization.
\end{enumerate}

D3 operates on the \emph{semantics} of the agent's output rather than the
\emph{form} of its input, making it robust to paraphrasing attacks that defeat
D1 and D2. However, it requires a precise definition of ``in-scope,'' which varies
across tasks, and may produce false positives on legitimate scope expansions.

\subsection{Defense Summary}

\begin{table}[t]
\centering
\caption{Properties of the three defenses.}
\label{tab:defenses}
\small\setlength{\tabcolsep}{5pt}
\resizebox{\columnwidth}{!}{%
\begin{tabular}{lccc}
\toprule
\textbf{Defense} & \textbf{+LLM calls} & \textbf{FP risk} & \textbf{Para.\ robust} \\
\midrule
D1: Provenance Tag   & 0  & Low    & Partial \\
D2: Struct.\ Extract & +1 & Medium & Yes \\
D3: Plan-Diff Gate   & +1 & High   & Yes \\
\bottomrule
\end{tabular}}
\end{table}
